\paragraph{名目総消費水準目標（NCLT）}
\label{sec:chapter5_beta_shock_policies_nclt}
式
\eqref{form:chapter3_policies_monetary_home_nclt_i_H_eq}、
\eqref{form:chapter3_policies_monetary_home_nclt_gamma_H_eq}および
\eqref{form:chapter3_policies_monetary_home_nclt_chi_H_eq}
により記述されたとおり、NCLTの式は以下のようであった。
\begin{align}
\label{form:chapter5_beta_shock_policies_nclt_i_H_eq}
    i_t^H&=i_{ss}^H+\phi_{gap}^H(\log(p_t^{H\to W}c_t^{H\to W})-\log(\chi_t^H))+\phi_{level}^H\gamma_t^H \\
\label{form:chapter5_beta_shock_policies_nclt_gamma_H_eq}
    \gamma_t^H&=\gamma_{t-1}^H+(\log(p_{t-1}^{H\to W}c_{t-1}^{H\to W})-\log(\chi_{t-1}^H)) \\
\label{form:chapter5_beta_shock_policies_nclt_chi_H_eq}
    \log(\chi_t^H)-\log(\chi_{ss}^H)&=\rho_{\chi}^H\left(\log(\chi_{t-1}^H)-\log(\chi_{ss}^H)\right)+\varepsilon_t^{\chi,H}
\end{align}
ここで
\(\phi_{gap}^H, \phi_{level}^H\)は反応係数、
\(p_t^{H\to W}c_t^{H\to W}\)は名目総消費水準、
\(\chi_t^H\)は目標パス、
\(\gamma_t^H\)は過去の乖離の累積、
\(\varepsilon_t^{\chi,H}\)は外生ショック項である。
\begin{enumerate}
    \item
    \label{itm:chapter5_beta_shock_policies_nclt_c_H_W}
        \(c_t^{H\to W}\)：
        総消費指数\(c_t^{H\to W}\)の図
        \ref{fig:chapter5_beta_shock_graphs_c_H_W}
        見ると、NCLTは100期までの前半と100期以降の後半の双方において好成績を収めている。
        そのことをグラフに特徴のある他の政策と比較することで理解する。
        まず総消費指数\(c_t^{H\to W}\)そのものを目標とする総消費水準目標（CLT）は
        「前半には他の政策よりも大きいが、後半には他の政策よりも小さい」というグラフを描く。
        これは前半においては落ち込みを最も小さく抑え込めるものの、
        後半に入るとやがて目標としていた初期定常状態の水準へ戻ってしまうためである。
        これに対しNCLTは
        「前半にはCLTに次いで大きいが、後半にはCLTに次いで小さい」
        という特徴をもつ。
        後半ではCLTが初期定常状態の水準に復帰するのに対し、
        NCLTはそれよりわずかに高い正の値を長期にわたって維持し続ける。
        一方で、物価目標を掲げる消費者物価水準目標（CPLT）や消費者物価インフレ目標（IT-CPI）は
        「前半には他の政策よりも小さいが、後半には他の政策よりも大きい」
        という対照的な動きを見せる。
        これらの政策は後半に大きなオーバーシュートの山を作るため後半の評価は良いが、
        前半には総消費指数\(c_t^{H\to W}\)を大暴落させてしまう重大な欠点がある。
        NCLTはCLTが持つ前半の強さと、CPLTやIT-CPIが持つ後半の強さを併せ持っている。
    \item 
    \label{itm:chapter5_beta_shock_policies_nclt_y_H}
        \(y_t^H\)：
        生産\(y_t^H\)の図
        \ref{fig:chapter5_beta_shock_graphs_y_H}
        を見ると、生産\(y_t^H\)は総消費指数\(c_t^{H\to W}\)と類似の動きをしており、
        その特徴も同様に2つの政策群との対比によって説明できる。
        まず総消費水準目標（CLT）の生産は
        「前半には他の政策よりも多いが、後半には他の政策よりも少ない」
        という波形を描く。
        これに対しNCLTの生産は
        「前半にはCLTに次いで多いが、後半にはCLTに次いで少ない」
        というグラフになっている。
        一方で、物価目標を掲げる消費者物価水準目標（CPLT）や消費者物価インフレ目標（IT-CPI）は
        「前半には他の政策よりも少ないが、後半には他の政策よりも多い」
        という対照的な波形を描く。
    \item 
    \label{itm:chapter5_beta_shock_policies_nclt_pi_H}
        \(\pi_t^H\)：
        グロス・インフレ率\(\pi_t^H\)の図
        \ref{fig:chapter5_beta_shock_graphs_pi_H}
        を見ると、全体として定常状態からの乖離が最も小さい政策は
        生産者物価水準目標（PPLT）および生産者物価インフレ目標（IT-PPI）であることがわかる。
        100期までの前半においてNCLTはPPLTおよびIT-PPIに次いで乖離が小さい。
        100期以降の後半においてもNCLTはPPLTおよびIT-PPIに次ぐ
        名目GDP水準目標（NGDPLT）や生産者物価テイラールールと同程度の乖離の小ささを達成している。
        また後半の特徴として、生産水準目標（CLT）の乖離が大きい水準で安定的にとどまっていることが見て取れる。
    \item
    \label{itm:chapter5_beta_shock_policies_nclt_Delta_t_minus_1_H}
        \(\Delta_{t-1}^H\)：
        価格分散\(\Delta_{t-1}^H\)の図
        \ref{fig:chapter5_beta_shock_welfare_Delta_H}
        を見ると、全体として定常状態からの乖離が最も小さい政策は生産者物価水準目標（PPLT）であり、
        次いで生産者物価インフレ目標（IT-PPI）および名目GDP水準目標（NGDPLT）であることがわかる。
        NCLTは全体としてこれらの政策に次いで定常状態からの乖離を小さくとどめているが、
        特に前半ではIT-PPIよりも小さい乖離を達成している。
        全体として特徴的であることは、
        総消費水準目標（CLT）の100期以降の後半における乖離が著しく増大し続けていることである。
        また消費者物価テイラールール（TR-CPI）、消費者物価水準目標（CPLT）および
        消費者物価インフレ目標（IT-CPI）も70期から150期にかけて大きな山の頂点を作っていることが確認できる。
\end{enumerate}
以上の分析から、NCLTがなぜ最高評価を得たのかが明らかになる。
NCLTは、素早い初動による消費の下支えを維持しつつも、
名目支出という枠組みを通じて「労働供給の過熱抑制」および「価格分散コストの抑制」を同時に達成しているのである。
この実質と名目のハイブリッドな安定化能力こそが、厚生を最大化させた本質的な要因である。